\documentclass[9pt, leqno]{amsart}
\usepackage[center]{titlesec}
\usepackage{lipsum}

\usepackage{amsmath, amsthm, amssymb, amsfonts, mathrsfs,relsize, thmtools, mathtools, stmaryrd, mlmodern, graphicx, setspace, geometry, float, quiver, xspace, hyperref, framed, tcolorbox, mathrsfs, enumerate, etoolbox}
\usepackage[T1]{fontenc}
\usepackage[all]{xy}
\usepackage[utf8]{inputenc}
\usepackage[english]{babel}
\usepackage[dvipsnames]{xcolor}
\usepackage[document]{ragged2e}
\usepackage[center]{titlesec}

\xyoption{curve}
\xyoption{2cell}
\UseAllTwocells

%Theoremstyle

\theoremstyle{definition}
\newtheorem{defn}{Definition}[section]
\newtheorem{example}[defn]{Example}
\newtheorem{construction}[defn]{Construction}
\newtheorem{notation}[defn]{Notation}

\theoremstyle{remark}
\newtheorem{remark}[defn]{Remark}
\theoremstyle{plain}
\newtheorem{thm}[defn]{Theorem}
\newtheorem{lemma}[defn]{Lemma}
\newtheorem{prop}[defn]{Proposition}
\newtheorem{cor}[defn]{Corollary}

% other symbols
\newcommand{\isom}{\cong}
\DeclareMathOperator{\idem}{E}
\DeclareMathOperator{\supp}{supp}
\DeclareMathOperator{\End}{End}
\DeclareMathOperator{\Inc}{Inc}
\DeclareMathOperator{\img}{im}
\DeclareMathOperator{\jac}{J}
\DeclareMathOperator{\Fib}{\mathbf{Fib}}
\DeclareMathOperator{\fib}{\textit{P} \colon \mathcal{E} \rightarrow \mathcal{B}}
\DeclareMathOperator{\opfib}{\textit{P} \colon \mathcal{E}^{op} \longrightarrow \mathcal{B}^{op}}
\DeclareMathOperator{\comp}{\circ}
\DeclareMathOperator{\is}{\equiv}
\DeclareMathOperator{\forw}{\Longrightarrow}
\DeclareMathOperator{\back}{\Longleftarrow}
\DeclareMathOperator{\adj}{\dashv}
\DeclareMathOperator{\2cell}{\Longrightarrow}


\newcommand{\dfib}[3]{\cat{#1} \colon \mathbb{#2} \rightarrow \mathbb{#3}}
\newcommand{\compdfib}[4]{\cat{#1}_{#4} \colon \mathbb{#2}_{#4} \longrightarrow \mathbb{#3}_{#4}}
\newcommand*{\dcat}[1]{\mathbb{#1}}
\newcommand*{\bspan}[1]{\mathbb{S}\text{pan}(#1)}
\newcommand*{\slice}[1]{\mathcal{#1}^2}
\newcommand*{\slicec}[1]{\mathbf{Fib}/#1}
\newcommand*{\dspan}[1]{\mathbb{S}\mathbf{pan}(#1)}
\newcommand{\cat}[1]{\mathcal{#1}}
\newcommand{\Cat}{\mathbf{Cat}}
\newcommand{\cart}[2]{\bar{#1}(#2) \colon #1 ^\ast #2 \rightarrow #2}
\newcommand{\hspan}[5]{\xymatrix@=.7cm{#1 & #2 \ar[l]_{#4}  \ar[r]^{#5} & #3}}
\newcommand{\cospan}[5]{\xymatrix@=.7cm{#1 \ar[r]^{#4} & #2 & #3 \ar[l]_{#5} & }}
\newcommand{\mor}[3]{#1 \colon #2 \rightarrow #3}
\newcommand{\proar}[2]{\xymatrix@=.7cm{#1 \ar[r]|-{\shortmid} & #2}}
\newcommand{\inproar}[3]{\xymatrix@=.7cm{#2 \ar[r]|-{\shortmid}^{#1} & #3}}
\newcommand{\lproar}[2]{\xymatrix@=.7cm{#1 \ar[l]|-{\shortmid} & #2}}
\newcommand{\pushoutcorner}[1][dr]{\save*!/#1+1.2pc/#1:(1,-1)@^{|-}\restore}
\newcommand{\pullbackcorner}[1][dr]{\save*!/#1-1.2pc/#1:(-1,1)@^{|-}\restore}
\newcommand{\chapquote}[2]{\begin{quotation} \textit{#1} \end{quotation} \begin{flushright} - #2 \end{flushright} }
\newcommand{\cleav}[2]{\varphi(#1,#2)}
\newcommand{\ter}[1]{\mathbf{1}_{\mathcal{#1}}}

\def\ie{i.e.\@\xspace}

%\usepackage[document]{ragged2e}

%\makeatletter
%\renewenvironment{abstract}{%
%	\begin{center}%
%		\scshape \abstractname% center the word "Abstract" in small caps (optional)
%	\end{center}%
%	\quotation
%	\noindent\centering
%}{%
%	\endquotation
%}
%\makeatother

%\title{Fibred Tabulators}
%\author{Siyabonga Mthimkhulu}
%\date{\today}
\titleformat{\section}
{\normalfont\Large\bfseries}
{\thesection}
{1em}
{\Large}

\titleformat{\subsection}
{\normalfont\large\bfseries}
{\thesubsection}
{1em}
{}

\titleformat{\subsubsection}
{\normalfont\large\bfseries}
{\thesubsubsection}
{1em}
{}


\begin{document}
	\title{Fibred Tabulators} %%%%%%%%%%%%
	\author[S.P Mthimkulu]{S.P Mthimkulu}
	\date{\today}
	\address{5 Melrose Road, Muizenberg, Cape Town, South Africa}
	\email{siyam@aims.ac.za}
	\maketitle
	

\begin{abstract}
	
	This paper defines tabulators in strict double fibration (\ie, double fibrations of double categories with tabulators). Additionally, we give a Beck-Chevalley Condition for fibred tabulators. Moreover, we show that  fibred tabulators are characterized by a preservation property.
\end{abstract}
	
\let\thefootnote\relax
\footnotetext{MSC2020: Primary 00A05, Secondary 00A66.} %%%%%%%%%%

\bigskip

\section*{Introduction}
\lipsum[2][1-4]

\section{Fibrations}

In this section we recall some basic facts about fibrations and their examples.

\begin{defn}
	Let $\fib$ be a functor. A morphism $\mor{f}{X}{Y}$ in $\cat{E}$ is \emph{cartesian} over $\mor{u}{I}{J}$ \ie, $P(f) = u$ if for any $\mor{v}{K}{I}$ in $\cat{B}$ and $\mor{h}{Z}{Y}$ over $u \comp v$ there is a unique morphism $\mor{g}{Z}{X}$ over $v$ with $h = f \comp g$.
	\begin{equation*}
		\xymatrix{
			\cat{E} \ar[d]_{P} && Z \ar@{-->}[r]_{g} \ar@/^1.3pc/[rr]^{h} & X \ar[r]_{f} & Y\\
			\cat{B} && K \ar[r]_{v} & I \ar[r]_{u} & J
		}
	\end{equation*}
	
	A morphism $\mor{k}{X}{Y}$ is called \emph{vertical} if it lies over an identity in $\cat{B}$. For an object $J \in \cat{B}$, the fibre of $P$ at $J$ denoted by $P(J)$ is the subcategory of $\cat{E}$ whose morphisms $k$ lie over $id_J$.
\end{defn}

\begin{defn}
	A functor $\fib$ is called \emph{fibred category} or a \emph{fibration} if for every $Y \in \cat{E}$ and $\mor{u}{I}{PY}$ there is a cartesian morphism $\cart{u}{E}$ over $u$. If $P$ is a fibration, a \emph{cleavage} for $P$ denoted $\varphi(u,E)$ is a choice of cartesian lifts for $u$ at $E$. Whence $P$ is called a cloven fibration.
\end{defn}

\subsection{Examples}

Let $\cat{C}^2$ denote the slice category of $\cat{C}$ over itself
\begin{example}[Fundamental Fibration]
	Let $\cat{C}$ be a category with pullbacks. Then the functor $\mor{\text{cod}}{\cat{C}^2}{\cat{C}}$ is a fibration; a commutative square is a cartesian arrow in $\cat{C}^2$ precisely when it is a pullback.
\end{example}

\begin{example}
	The domain functor $\mor{\text{dom}}{\slice{\cat{C}}}{\cat{C}}$ is a fibration.
\end{example}

\begin{example}[Family fibration]
	Every category $\cat{C}$ gives rise to a \emph{family fibration} $\mor{\cat{F}}{Fam(\cat{C})}{Set}$. A morphism $(u,\{\varphi_i\}_{i \in I})$ is cartesian when every $\varphi_i$ is an isomorphism. Then $\cat{F}$ is a fibration since given $\mor{u}{I}{J}$ and $\{Y_j\}_{j \in J}$, $\mor{(u,\{1_{Y_{u(i)}}\})}{\{Y_{u(i)}\}_{i \in I}}{\{Y_j\}_{j \in J}}$ is cartesian over $u$. 
\end{example}

\subsection{Fibred Finite Limits}

In \cite{TS} fibred finite limits are given in terms of terminal objects and pullbacks. Furthermore, \cite{JacobsCLTT} captures the notion of fibred terminal objects in $\mathbf{Fib}$ using fibred adjunctions. 

\begin{defn}
	A fibration $\fib$ has fibred terminal objects if every fibre $P(I)$ has a terminal object $1_I$ and they are preserved by cartesianness.
\end{defn}

In $\slicec{\cat{B}}$, the category of fibrations over $\cat{B}$, a fibration $P$ has fibred terminals if the fibred functor $\mor{F}{P}{id_B}$ has a right adjoint. Moreover, this property is characterized by preservation as follows.

\begin{lemma}\label{fibredterminals}
	Let $\cat{B}$ be a category with a terminal object $\ter{B}$. Then $\fib$ has fibred terminals if and only if $\cat{E}$ has a terminal $1_\cat{E}$ and $P(\ter{E})$ is a terminal in $\cat{B}$.
\end{lemma}

\begin{proof}
	Consider the following diagram in $\Fib$
	
	\begin{equation*}
		\xymatrix{
			\cat{E} \ar[dd]_{P} \ar@<1.4ex>[rr]^{!} & \perp & 1 \ar@<1.4ex>[ll]^{\ter{E}} \ar@{=}[dd]^{id}\\
			\\
			B \ar@<1.4ex>[rr]^{!} & \perp & 1 \ar@<1.4ex>[ll]^{\ter{B}}
		}
	\end{equation*}
\end{proof}

Futhermore, we have the following theorem

\begin{thm}
	If $\cat{B}$ has a terminal object $\ter{B}$ and $\fib$ is a fibration, the following are equivalent
	\begin{enumerate}
		\item $P$ has a terminal object in $\slicec{\cat{B}}$. \\
		\item  $P$ has a terminal object in $\Fib$.
	\end{enumerate}
\end{thm}

\begin{proof}
	By the lifting and factorization of adjoints
	
	\begin{equation*}
		\xymatrix{
			\cat{E} \ar[dr]_{P} \ar@/^1pc/[rr]^{!} \ar[r]^{P} & B \ar@{=}[d] \ar[r]^{!} & 1 \ar@{=}[d]^{id}\\
			& B \ar[r]_{!} & 1
		}
	\end{equation*}
	Given $!_{\cat{B}} \adj \ter{B}$, it follows that $P \adj \ter{E}$ if and only if $!_\cat{E} \adj \ter{E}$ over $!_\cat{B} \adj \ter{B}$.
\end{proof}

\begin{defn}
	A fibration $\fib$ has fibred pullbacks if every fibre has pullbacks and they are preserved by cartesianness.
\end{defn}

Similarly, fibred pullbacks can be characterised as follows.
\begin{thm}\label{fibredpb}
	Let $\cat{B}$ be a category with pullbacks. A fibration $\fib$ has fibred pullbacks if and only if $\cat{E}$ has and $P$ preserves pullbacks. 
\end{thm}

The following result is a consequence of Lemma \ref{fibredterminals} and Theorem \ref{fibredpb} 

\begin{thm}
	For a category $\cat{B}$ with finite limits, a fibration $\fib$ has fibred finite limits if and only if $\cat{E}$ has finite limits and $P$ preserves finite limits. 
\end{thm}



\bibliographystyle{alpha}
\bibliography{reference}


\end{document}